\documentclass[12pt,a4paper]{article}
\usepackage[utf8]{inputenc}
\usepackage[spanish,es-noquoting]{babel}
\usepackage[T1]{fontenc}
\usepackage{lmodern}
\usepackage{geometry}
\geometry{top=2cm,bottom=2cm,left=2.5cm,right=2.5cm}
\usepackage{xcolor}
\usepackage{listings}
\usepackage{enumitem}
\usepackage{fancyhdr}
\usepackage{graphicx}
\graphicspath{{../../assets/}}
\lstset{language=C,basicstyle=\ttfamily\small,breaklines=true,frame=single,showstringspaces=false}
\pagestyle{fancy}
\setlength{\headheight}{15pt}
\fancyhf{}
\lhead{Monitorías Sistemas Operativos 2026-II}
\rhead{Capítulo 4: Sincronización}
\cfoot{\thepage}
\title{\textbf{Taller Propuesto: Coordinación de Hilos}}
\author{Monitorías Sistemas Operativos}
\date{\today}
\begin{document}
\maketitle

\section*{Propósito}
Este taller trabaja exclusión mutua, barreras y variables de condición. Compile con \texttt{-pthread}.

\section*{Ejercicio 1}
Hay $C$ cupos disponibles en un laboratorio. Se crean $P$ hilos, cada uno con una solicitud de cupos.
\begin{enumerate}
    \item Proteja con un mutex la lectura, validación y actualización de los cupos.
    \item Una solicitud se aprueba completa o se rechaza; nunca deben quedar cupos negativos.
    \item Al finalizar, el \texttt{main} debe informar solicitudes aprobadas, rechazadas y cupos restantes.
    \item Explique qué intercalamiento podría producir una sobreasignación sin mutex.
\end{enumerate}

\section*{Ejercicio 2}
Un vector de $N$ enteros se divide entre cuatro hilos. En la primera fase cada hilo calcula la suma de su bloque. En la segunda, todos deben normalizar sus elementos restando el promedio global entero.
\begin{enumerate}
    \item Use una barrera para impedir que cualquier hilo inicie la segunda fase antes de que todos hayan producido su suma parcial.
    \item Designe un hilo para calcular el promedio global entre las dos fases.
    \item Use una segunda barrera para asegurar que los demás hilos no lean el promedio antes de que esté disponible.
\end{enumerate}

\section*{Ejercicio 3}
Implemente una cola circular de capacidad $K$ para solicitudes de impresión. Dos hilos productores generan documentos identificados y un hilo consumidor los atiende.
\begin{enumerate}
    \item Proteja la cola con un mutex y use dos variables de condición: \texttt{no\_llena} y \texttt{no\_vacia}.
    \item Cada productor genera exactamente $M$ documentos; el consumidor debe terminar luego de atender $2M$.
    \item Use \texttt{while} al esperar las condiciones y explique qué error aparecería con \texttt{if}.
\end{enumerate}

\section*{Ejercicio 4}
Cuatro hilos actualizan marcadores de equipos en un arreglo compartido. Un hilo observador debe imprimir el tablero solo cuando todos los trabajadores hayan concluido una ronda.
\begin{enumerate}
    \item Use un mutex para proteger cada actualización del arreglo.
    \item Use una barrera para delimitar el final de cada una de tres rondas.
    \item Asegure que el observador imprima una ronda completa y que los trabajadores no comiencen la siguiente antes de esa impresión.
\end{enumerate}
\end{document}
